\documentclass[12pt]{article}
\usepackage[margin=1in]{geometry}
\usepackage{booktabs}
\usepackage{siunitx}
\sisetup{round-mode=places,round-precision=3}

\begin{document}
\thispagestyle{empty}
\begin{center}
    {\LARGE Dissertation Project}
\end{center}
\newpage

\section{Results}
This section reports quantitative findings from implementation outputs in \texttt{reports/tables/}, with emphasis on the Phase 3 evaluation pipeline (\texttt{src/evaluation/}).

\subsection{Core Classification Metrics (\texttt{src/eval/metrics.py})}
The \texttt{src/eval/metrics.py} module computes the core multiclass performance metrics used as the baseline model-quality summary. These are exported via \texttt{overall\_metrics\_*.csv} and \texttt{per\_class\_metrics\_*.csv}.

For the selected ensemble (\texttt{overall\_metrics\_ensembleselected.csv}), results were:
\begin{itemize}
    \item Accuracy: 0.737
    \item Macro precision: 0.733
    \item Macro recall: 0.739
    \item Macro-F1: 0.735
    \item Weighted-F1: 0.733
    \item Macro AUC (OvR): 0.904
    \item Log-loss: 0.553
    \item Expected Calibration Error (ECE): 0.044
\end{itemize}

Per-class performance (\texttt{per\_class\_metrics\_ensembleselected.csv}) shows strongest discrimination for AD (AUC: 0.955) and SCD (AUC: 0.938), with MCI as the most difficult class (AUC: 0.818; F1: 0.605). This class-level asymmetry is clinically relevant because the intermediate stage (MCI) is the hardest boundary.

\subsection{Complete \texttt{src/eval/} Pipeline Outputs}
The full \texttt{src/eval/} stack is orchestrated by \texttt{src/eval/orchestrator.py}, which combines 11 components into one evaluation pass. In this project, it generated model-wise outputs for CatBoost, Random Forest, Neural Net, Ensemble, and EnsembleSelected.

\textbf{1) Metrics (\texttt{metrics.py})}
\begin{itemize}
    \item \texttt{reports/tables/overall\_metrics\_*.csv}
    \item \texttt{reports/tables/per\_class\_metrics\_*.csv}
\end{itemize}

\textbf{2) Visual diagnostics (\texttt{visualizations.py})}
\begin{itemize}
    \item \texttt{reports/figures/confusion\_matrix\_*.png}
    \item \texttt{reports/figures/roc\_curves\_*.png}
    \item \texttt{reports/figures/calibration\_curves\_*.png}
\end{itemize}

\textbf{3) Statistical inference (\texttt{statistical\_inference.py})}
\begin{itemize}
    \item \texttt{reports/tables/bootstrap\_ci\_*.csv}
\end{itemize}
These tables provide bootstrap means and confidence intervals for overall and per-class metrics, used to quantify stability.

\textbf{4) Uncertainty characterization (\texttt{uncertainty.py})}
\begin{itemize}
    \item Exported into consolidated reports/JSON summaries via orchestrator outputs
\end{itemize}
Uncertainty signals (entropy, margin, max probability) are used to identify difficult cases and support escalation logic.

\textbf{5) Error taxonomy (\texttt{error\_taxonomy.py})}
\begin{itemize}
    \item \texttt{reports/tables/error\_taxonomy\_*.csv}
    \item \texttt{reports/tables/hitl\_review\_*.csv}
\end{itemize}
This gives row-level error records and a shortlist of highest-uncertainty errors for clinician review.

\textbf{6) Explainability (\texttt{explainability.py})}
\begin{itemize}
    \item \texttt{reports/figures/shap\_summary\_catboost.png}
    \item \texttt{reports/tables/feature\_importance\_catboost.csv}
\end{itemize}
SHAP is enabled in this environment primarily for CatBoost, providing global feature attribution for interpretability.

\textbf{7) Bias diagnostics (\texttt{bias\_diagnostics.py})}
\begin{itemize}
    \item Group fairness outputs when demographic metadata are provided to orchestrator
\end{itemize}
This is complemented by the newer dedicated fairness module in \texttt{src/evaluation/fairness\_analysis.py}.

\textbf{8) Consolidated reporting (\texttt{reporting.py})}
\begin{itemize}
    \item \texttt{reports/reports/consolidated\_report\_*.csv}
    \item \texttt{reports/reports/results\_*.json}
\end{itemize}
These files aggregate all component outputs for each model and provide audit-ready machine-readable summaries.

\textbf{Interpretation:} the \texttt{src/eval/} folder provides the baseline technical audit layer (model quality, uncertainty, error forensics, explainability, and reporting), while \texttt{src/evaluation/} adds Phase 3 experimental analyses (HITL comparisons, threshold/cost/fairness/agreement/efficiency studies).

\subsection{AI-only vs HITL Comparison}
From \texttt{statistical\_tests.csv}, AI-only baseline accuracy was 0.737 and HITL accuracy was 0.828 (absolute gain: 0.091). Macro-F1 improved from 0.735 to 0.827 (gain: 0.092).

McNemar's paired test showed a statistically significant difference between AI-only and HITL performance ($\chi^2=36.54$, $p=1.49\times10^{-9}$), with matched odds ratio 22.00 (note: $b=2$, $c=44$).

Bootstrap confidence intervals (\texttt{statistical\_tests.csv}) support this improvement:
\begin{itemize}
    \item Accuracy difference (HITL - AI): 95\% CI [0.065, 0.120]
    \item Macro-F1 difference (HITL - AI): 95\% CI [0.065, 0.121]
\end{itemize}

From \texttt{hitl\_simulated\_clinician\_experiment.csv}, the active feedback experiment showed the following best performance by clinician accuracy level:

\begin{table}[h]
\centering
\caption{Best HITL result by simulated clinician accuracy}
\begin{tabular}{lccc}
\toprule
Clinician accuracy & Best round & Best test accuracy & Gain vs round 0 \\
\midrule
0.85 & 1 & 0.784 & +0.052 \\
0.90 & 4 & 0.802 & +0.069 \\
0.95 & 5 & 0.776 & +0.043 \\
\bottomrule
\end{tabular}
\end{table}

\subsection{Threshold Sensitivity Analysis}
From \texttt{threshold\_sensitivity\_summary.csv}, the selected threshold was stable at 0.69 across assumed clinician accuracies 0.85--0.95. Review rate remained 0.196, while policy accuracy increased from 0.808 to 0.827 as clinician accuracy increased.

This indicates the threshold choice is robust in the tested range and gives a consistent review burden with predictable accuracy gains over AI-only performance.

\subsection{Ablation Study on Escalation Signals}
From \texttt{ablation\_study.csv}, removing \texttt{risk\_2\_margin} produced the largest performance drop (accuracy change: -0.024; AUC change: -0.0076), suggesting this signal contributes most strongly among tested features.

Some removals had negligible effect (e.g., \texttt{risk\_1\_inv\_conf}), while removing \texttt{risk\_3\_entropy} slightly increased performance on this split. This suggests partial redundancy among signals and motivates future feature regularization/selection.

\subsection{Cost-Benefit Analysis}
From \texttt{cost\_analysis.csv}, HITL reduced clinical cost relative to AI-only across tested thresholds and human-accuracy settings. Best observed cost reductions were:

\begin{table}[h]
\centering
\caption{Best observed cost reduction by simulated clinician accuracy}
\begin{tabular}{lccc}
\toprule
Clinician accuracy & Threshold & Review rate & Cost reduction \\
\midrule
0.85 & 0.50 & 0.448 & 115.0 \\
0.90 & 0.28 & 0.724 & 137.0 \\
0.95 & 0.24 & 0.767 & 171.5 \\
\bottomrule
\end{tabular}
\end{table}

This supports the claim that HITL can improve clinically weighted utility, not only raw accuracy.

\subsection{Fairness Analysis}
From \texttt{fairness\_analysis.csv} and \texttt{fairness\_summary.csv}:
\begin{itemize}
    \item Overall accuracy: 0.737; overall escalation rate: 0.817.
    \item Gender disparity: demographic parity diff = 0.009, equalized odds diff = 0.087.
    \item Age-bin disparity: demographic parity diff = 0.163, equalized odds diff = 0.147.
\end{itemize}

The age-group disparities are materially larger than gender disparities, indicating age-related fairness risk in the current escalation behavior.

\subsection{Operational Efficiency and AI-Human Agreement}
From \texttt{efficiency\_summary.csv} and \texttt{agreement\_summary.csv}, currently logged dashboard data are minimal (1 review event): median decision time was 201.47 seconds, agree rate was 1.0, and disagree rate was 0.0.

These early values are not yet statistically stable and should be treated as instrumentation validation rather than final efficiency/agreement conclusions.

\subsection{Uncertainty Quantification}
From \texttt{conformal\_prediction.csv}, empirical conformal coverage was 1.0 for target coverages 0.80--0.95, with average prediction-set size between 2.34 and 2.64 classes. Singleton sets were rare (0.0--0.026), indicating conservative uncertainty sets.

From \texttt{uncertainty\_decomposition.csv}, mean total uncertainty was 0.677, mostly aleatoric (mean 0.651), with smaller epistemic uncertainty (mean 0.025; epistemic fraction mean 0.040).

\subsection{Summary of Evidence}
Across statistical, cost, ablation, fairness, and uncertainty analyses, the current implementation provides consistent evidence that the HITL pipeline improves predictive and clinical utility over AI-only baseline. The main remaining technical risks are (i) fairness disparity across age groups and (ii) limited logged interaction data for stable efficiency/agreement estimates.

\end{document}
